\section{Переобкачка}

Переобкачка вынесена в отдельную программу, \texttt{recrawler.py}. Она ничего не
скачивает \emph{в базу}: её задача --- определить, изменился ли документ, и если
да, поставить его URL в очередь. Дальше документ заберёт обычный робот.

Программа берёт из базы документы, которые давно не проверялись (порог задаётся
в секции \texttt{recrawl} конфига), и по каждому идёт в два шага --- от дешёвого
к дорогому.

\textbf{Шаг 1. Условный запрос.} Если у документа сохранены \texttt{ETag} или
\texttt{Last-Modified}, запрос отправляется с заголовками
\texttt{If-None-Match} и \texttt{If-Modified-Since}. Ответ \texttt{304 Not
Modified} означает, что документ точно не менялся, --- причём тело страницы
сервер при этом не передаёт вовсе.

\textbf{Шаг 2. Сравнение отпечатка.} Если сервер условные запросы не
поддерживает, страница скачивается, из неё выделяется текст и считается
отпечаток SHA-256. Он сравнивается с сохранённым в базе.

Сравнивается именно \emph{выделенный текст}, а не \enquote{сырой} HTML. Это
принципиально: в HTML на каждой перезагрузке меняются счётчики, идентификаторы
сессий, рекламные блоки и метки времени, и по нему почти любой документ выглядел
бы изменившимся. Отпечаток текста меняется только тогда, когда изменилось
содержание.

Изменившийся URL попадает в очередь со статусом \texttt{pending} и приоритетом
1. Приоритет важен: очередь отдаёт такие адреса раньше новых, поэтому обновление
уже известных документов не ждёт, пока разгребётся весь фронт обкачки. При
повторной обкачке запись в \texttt{documents} обновляется, а счётчик
\texttt{revision} увеличивается на единицу.

\subsection*{Проверка работы}

Цикл проверялся целиком. Сначала переобкачка запускалась на свежескачанных
документах --- все десять вернули \texttt{304 Not Modified}, в очередь не попал
никто. Затем у двух документов в базе отпечаток был подменён вручную, что
имитирует изменение страницы на сайте. Повторный запуск нашёл ровно эти два
документа и поставил их в очередь:

\begin{verbatim}
[изменился] https://ru.wikipedia.org/wiki/?_(фильм)
[изменился] https://ru.wikipedia.org/wiki/…а_пятый_всадник_—_Страх

Проверено документов:      6
  ответ 304 Not Modified:  4
  текст совпал:            0
  изменились:              2
В очередь на переобкачку добавлено: 2
\end{verbatim}

Запущенный следом робот взял оба адреса первыми, до новых документов, и поднял
им \texttt{revision} с 1 до 2:

\begin{verbatim}
[updated ]   15790 симв.  https://ru.wikipedia.org/wiki/?_(фильм)
[updated ]    3011 симв.  https://ru.wikipedia.org/wiki/…а_пятый_всадник_—_Страх
[new     ]    1863 симв.  https://ru.wikipedia.org/wiki/0-41*
\end{verbatim}

\pagebreak
